\chapter{当代因果哲学：反事实、概率、机制与干预}
\label{ch:philosophies}

\begin{learninggoals}
读完本章，你应当能够：比较反事实、概率、机制、能力与干预理论的分析对象；检验先占、过度决定、遗漏和共同原因等反例；说明多元证据怎样形成受约束的整合，而非概念拼贴。
\end{learninggoals}

\section{反事实因果：刘易斯及其难题}

\subsection{反事实语义、依赖与因果链}

日常判断原因时，我们自然会问：“如果原因没有发生，结果还会发生吗？”若苏茜投出的石块击碎玻璃，而没有这次投掷玻璃便不会碎，我们倾向于把投掷视为原因。这个\term{若无检验}（\eng{but-for test}）抓住了原因作为差异制造者的直觉，并给出因果方向：结果依赖原因，原因通常不以同样方式依赖结果。
困难在于反事实前件与事实相反。我们不能简单查看实际数据中“同一个对象同时未受处理”的结果，也不能只用材料蕴涵，因为一个前件为假的材料条件句会自动为真。反事实理论需要说明：在保留多少实际背景的前提下，怎样评价“若 $A$ 未发生，$C$ 会怎样”。

\paragraph{可能世界与相似性}

\textcite{lewisCounterfactuals}以可能世界之间的比较相似性给出反事实语义。常用的直观说法是：考察与现实最相似的 $A$ 世界，若相关的最相似 $A$ 世界都是 $C$ 世界，则“若 $A$，则 $C$”为真。但Lewis并不要求必须存在一个或一组“最近的 $A$ 世界”；他明确拒绝这种\term{极限假设}（limit assumption）。在不假定最近世界存在的更一般表述中，只要存在 $A$ 世界，“若 $A$，则 $C$”为真可理解为：有某个 $A\land C$ 世界比任何 $A\land\neg C$ 世界都更接近现实。因此，“最接近世界”适合作直观入口，却不应被读成Lewis正式语义必须满足的存在性条件。
这里的世界不是平行宇宙的经验假说，而是评价模态命题的理论装置。世界之间的“接近”也不是表面相似度的简单加总，而由受语境与自然规律约束的比较相似性排序决定。

评价“若火柴没有被划，是否会燃烧”时，我们通常允许一个小的局部偏离使划火柴事件不发生，同时尽可能保留现实规律和广泛事实。若为了保持更多局部事实而允许大规模违反规律，反事实判断会失去稳定性。Lewis在讨论因果时给予避免巨大、广泛奇迹较高优先级，再考虑事实吻合范围。
相似性语义的优点是能够处理非因果反事实，并解释为何“回溯式”反事实通常不成立。例如，若结果没有发生，我们通常不会反推原因也未发生，而倾向于假定因果过程在原因之后出现局部偏离。但相似性排序本身也受到批评：若我们已经凭因果判断决定哪些世界更接近，再用接近关系定义因果，就可能产生循环。

\paragraph{从反事实依赖到因果链}

Lewis早期分析把事件理解为时空区域中具有特定性质的个体。若 $c$ 和 $e$ 都实际发生，且
\[
\neg O(c)\xRightarrow{\mathrm{cf}} \neg O(e),
\]
即在最接近的无 $c$ 世界中也没有 $e$，则 $e$ 对 $c$ 存在\term{因果依赖}\parencite{lewisCausation1973}。这里 $O(c)$ 表示事件 $c$ 发生。

简单依赖并不总是传递。若按下开关导致继电器动作，继电器动作导致灯亮，我们希望开关是灯亮的远因，即使某些复杂背景使端点之间的直接反事实依赖失败。Lewis因此把因果关系定义为因果依赖链的祖先关系：存在事件序列 $c,d_1,\dots,e$，每个后项因果依赖前项。
这一做法恢复了因果传递性，却也制造过度传递问题。一个人向右躲开巨石，因此恰好站到后来被树枝击中的位置；躲避依赖巨石，受伤依赖躲避，但说巨石造成树枝击伤可能很勉强。是否传递取决于我们追踪的是能量过程、机会条件，还是与问题相关的变量路径。

\subsection{先占、过度决定与遗漏难题}

\term{早期先占}的经典例子中，苏茜和比利都向瓶子投石。苏茜的石块先击中并打碎瓶子，比利的石块稍后到达原位置。苏茜是实际原因，比利只是备份。然而，若苏茜没有投掷，比利仍会打碎瓶子，所以“没有苏茜投掷，瓶子不会碎”为假。简单若无检验遗漏了苏茜这个原因。

因果链方法可以区分实际过程：苏茜投掷经由飞行、撞击和裂纹形成到达破碎；比利链条在撞击前被截断。但\term{晚期先占}更困难：两条过程可以延伸很远，直到后期一个过程阻止另一个。要指出“哪条链是实际链”，理论往往需要对过程结构作比简单端点反事实更细的描述。

\term{对称过度决定}中，两枚同时且各自足够的炸弹共同摧毁目标。没有任意一枚，另一枚仍足够，因而两者都不通过若无检验；直觉却把两者都视为原因。可以把检验改成“在适当背景或某个事件子集中，改变该因素会改变结果”，但何谓适当背景必须防止任意挑选。

\paragraph{从是否发生到如何发生}

Lewis后期提出\term{因果影响}分析。一个事件 $c$ 影响事件 $e$，不是只要求“无 $c$ 则无 $e$”，而是要求 $c$ 的一系列细微变化与 $e$ 的一系列变化之间存在系统反事实对应\parencite{lewisInfluence2000}。投石时间、速度或方向的变化会改变玻璃破碎的时间、方式或碎片分布；即使备份原因保证“最终仍破碎”，实际结果的具体方式仍受苏茜投掷影响。
这一修正把二值发生扩展为事件的时间和方式，使先占案例更可处理。它仍依赖恰当的事件变更和相似性排序，也可能在微小改变引发复杂连锁时过宽。值得注意的是，《作为影响的因果》发表于2000年，而非有些二手资料误记的1986年。

\paragraph{遗漏、双重防止与规范背景}

“园丁没有浇水导致植物死亡”把一个缺席当作原因。若事件是时空中的具体个体，缺席似乎不是可进入因果链的事件；若允许一切缺席成为原因，则任何未发生的救援都可能被说成结果原因。我们通常依赖职责、常规或期待筛选相关遗漏：受雇园丁的不浇水比陌生人的不浇水更容易被视为原因。

\term{双重防止}也挑战简单过程直觉：保镖阻止刺客阻止解药送达，因而保镖的行动使解药到达。反事实上，保镖行动与获救存在依赖，但两者之间未必有从原因向结果传递的连续物理过程。反事实理论能自然容纳这类“让某过程不受阻”的因果，纯能量传递理论则较困难。

这些案例表明，实际因果判断可能同时依赖反事实结构、过程连接、事件描述粒度与正常性规范。结构因果模型中的实际因果定义尝试把这些要素形式化，但不同定义在先占、投票、法律责任等案例上仍有分歧。

\subsection{因果类型、事件粒度与证据来源}

“吸烟导致肺癌”通常是\term{类型因果}命题，涉及人群中的暴露分布与风险变化；“张某的吸烟导致了他的肺癌”是\term{个例因果}命题。前者可由试验、观测设计和机制证据支持，后者还需评价该个体的替代历史与实际过程。群体平均效应为零并不排除对不同亚组存在方向相反的作用，平均效应为正也不自动确定每一患者的实际原因。
现代政策评估大多以类型因果效应为目标；法律、历史和故障诊断则频繁提出个例因果问题。选用框架时应先看查询，而不是假定所有因果句都由同一个若无条件定义。

\paragraph{深读案例：委员会投票与结构敏感性}

五人委员会以至少三票通过提案，实际五人都投赞成。对任一成员，若只有他改投反对，提案仍通过，所以简单若无检验不把任何一票视为原因。若允许把其他两票设为反对作为应急背景，则某一票会成为关键；同样推理又把每票都算作原因。
这是否正确取决于查询。若问每位成员是否参与形成充分多数，全部赞成票都有贡献；若问谁的改变本可阻止提案，单独没有人；若问法律责任，职务、信息和串谋可能改变答案。结构模型需要明确聚合规则和应急情形，不能用一个“实际原因”标签消除问题差异。
再设主席只在平票时投票，实际普通成员三比一通过，主席未投票。主席的缺席不是提案通过的原因；但制度规则允许他的潜在票在其他背景中关键。这里必须区分因果能力、实际贡献和制度权力。反事实理论若不限定实际过程，可能把广泛备份能力误作具体原因。

\paragraph{事件描述粒度}

“苏茜投石”可以描述为一次具体肌肉动作、石块以某速度离手、石块击中某位置，或更抽象的“有人投掷”。反事实依赖会随描述变化。若苏茜没有以每秒十米速度投掷，她可能以九米速度投掷，瓶子仍碎；若她完全没有投掷，比利会接替。粗细粒度决定哪些替代世界最接近。
结果描述也一样。“瓶子碎了”在备份存在时不依赖苏茜，但“瓶子在12点00分以这种裂纹碎裂”可能依赖。后期影响理论利用这种细节，但过细结果会使几乎任何微小背景都成为原因。适当粒度应由解释问题决定，并保持在相似案例中一致，而不是为了得到直觉答案临时调整。
结构方程把事件粒度转化为变量选择。二元变量适合某些责任规则，过程变量适合先占路径。模型比较应说明新增变量如何对应真实过程及可获得证据，而非只证明“存在一幅图能给出想要答案”。

\paragraph{反事实判断的证据来源}

我们无法直接观察反事实，却可由多类证据约束。重复试验估计类型层面的干预差异；机制知识说明改变如何传播；时间与物理连续性识别实际路径；相似案例支持背景稳定；制度规则规定哪些替代情形可行。个例因果判断通常综合这些来源。
例如判断药物是否导致某患者过敏，群体试验显示风险提升，免疫检测显示特异机制，给药后时间吻合，停药后症状缓解，再暴露后复发。任何单项证据都可能不足，组合则强化“若无给药便无此次反应”的模型。反事实不是纯粹想象，而是受实际证据约束的模型推论。

\section{概率、机制、力量与干预：多元因果理论}

\subsection{概率、共同原因与INUS条件}

反事实理论突出差异制造，却可能难以刻画实际过程；规律论说明事件类型的稳定联系，却容易遗漏方向；机制论展示结果如何产生，却未必给出效应大小；概率理论适合不确定系统，却必须处理混杂；干预主义连接实验实践，却需要说明何种变化算合格干预。二十世纪后半叶的因果哲学由此呈现多元格局。
多元不等于所有理论都同样正确，也不意味着同一句因果话语可以任意解释。较稳妥的看法是：科学中的因果概念承担若干彼此相关的功能------预测操纵结果、解释生成过程、支持迁移、归属责任。不同理论把其中某一种功能当作分析中心。评价它们时，应同时问理论处理何种对象、服务何种查询、依赖何种背景。

\paragraph{概率提升与共同原因}

在随机或不完全观测系统中，原因不必保证结果。吸烟者并非人人患肺癌，未吸烟者也可能患病。概率因果理论据此把原因理解为提高结果概率的因素。最简单形式是
\[
P(E\mid C)>P(E\mid \neg C).
\]
Suppes把时间优先和概率提升结合起来讨论“表面原因”与“真正原因”\parencite{suppes1970}。

简单概率提升面临三类困难。第一是\term{混杂}：冰淇淋销售与溺水同时上升，是炎热天气共同影响二者，并非购买冰淇淋使人溺水。第二是\term{辛普森反转}：总体概率提升可能在每个相关亚组中消失或反向。第三是\term{预防原因}：安全带对死亡具有负向因果作用，表现为降低而非提高概率。因而更一般的语言应谈论概率改变，并要求在恰当背景中比较。

Reichenbach的\term{共同原因原则}指出：若两个事件相关，且彼此不是原因，则应寻找一个共同原因，使相关在适当条件下被“筛除”\parencite{reichenbach1956}。在简单情形中，若 $C$ 是 $A$ 与 $B$ 的共同原因，可期待
\[
P(A,B\mid C)=P(A\mid C)P(B\mid C).
\]
现代图模型把这种思想推广为条件独立结构，但共同原因原则不是无条件定理。选择偏差、确定关系、量子相关和不完整变量描述都会使简单筛除形式需要修正。

\paragraph{Mackie的INUS条件}

现实结果常由原因复合产生。电路短路并非任何情况下都引起火灾；它可能需要可燃物、氧气和保护装置失效共同出现。Mackie把许多日常原因分析为\term{INUS条件}：一个不充分但为某个不必要而充分条件之必要部分\parencite{mackie1974}。

设火灾 $F$ 可以由两组条件产生：
\[
(S\land O\land D)\lor(L\land G)\longrightarrow F,
\]
其中 $S$ 为短路，$O$ 为有氧环境，$D$ 为保护失效；另一条路径由雷击 $L$ 与干燥 $G$ 构成。$S$ 单独不足以导致火灾，但在组合 $S\land O\land D$ 中不可缺；这个组合足以导致火灾，却不是必要的，因为还有另一条路径。因此 $S$ 是 $F$ 的一个INUS条件。

INUS分析有两项贡献。它避免要求单一原因既必要又充分，并突出因果场景与替代充分集合。但形式中的“充分”通常相对于背景和模型，而非宇宙中绝对无例外；如何选择正常背景、如何处理概率结果和连续变量，也不是析取范式本身能够解决。INUS更像揭示复合原因结构的分析工具，而不是完整的因果形而上学。

\paragraph{Salmon：解释必须进入因果网络}

Salmon反对把科学解释仅仅看成定律下的推导。他早期使用统计相关模型，后来发展\term{因果过程}理论：能够传递结构或“标记”的过程，与只呈现运动轨迹的伪过程不同\parencite{salmon1984}。一束光可以携带调制信息；墙上的光斑可以超光速移动，却没有一个物体或信号沿光斑轨迹传播。
因果过程在交互中改变并把变化继续传递。两辆汽车碰撞后速度和形变改变，交互把此前过程连接到此后过程。解释一个事件，就是说明它如何嵌入产生它的因果网络，而不是仅仅展示一个包含它的统计规律。
“标记传递”后来受到批评：什么算可传递标记可能依赖反事实，难以覆盖场、持续约束和非局部结构。Salmon转向以守恒量传递刻画因果过程，但这更适合物理情形，对社会制度和生物调控未必直接适用。其核心遗产不是某个永恒定义，而是解释相关性要求追踪实际产生结构。

\subsection{机制、因果力量与干预主义}

生命科学中的机制解释通常不要求一条单独守恒量链，而把机制理解为实体和活动以特定时空方式组织并产生现象。Machamer、Darden与Craver强调“实体做什么”，Glennan则强调部件相互作用形成稳定行为\parencite{mdc2000,glennan2017}。
以动作电位为例，完整机制涉及膜、电压门控离子通道、离子浓度梯度以及通道依时间开闭的组织。仅说刺激与放电相关不能解释波形如何产生；仅列出部件也不够，必须说明活动顺序和组织。机制模型往往是多层且不完整的：研究目的决定需要向下分解到何种程度。
机制证据对外推尤其重要。若某药物在试验人群中有效，而目标人群缺少关键受体，平均效应未必可迁移。不过，机制叙述也可能只是“故事”。可靠机制解释需要独立测量、干预部件、追踪中间过程并排除竞争路径。

\paragraph{因果力量与能力}

Cartwright批评把自然定律视为处处无条件成立的普遍概括。实验定律常描述精心隔离的情形；现实结果由多种\term{因果能力}共同作用\parencite{cartwright1989}。阿司匹林具有缓解疼痛的能力，并不意味着每次服用都实际止痛，因为剂量、吸收、病因和其他药物会改变表现。
力量或能力语言抓住了原因在不同背景中仍可被识别的倾向性，也说明效应不是裸相关。但它面临测量和判准问题：若我们只通过效果知道某物有能力，再用能力解释效果，便有循环危险。实际科学需要通过多种独立表现、机制和干预来校准能力归属。

\paragraph{Woodward：不变性与干预}

Woodward的\term{干预主义}不把因果简化为人类实际能够操纵的东西，而是用理想干预定义变量层面的因果关系\parencite{woodward2003}。粗略地说，若存在对 $X$ 的合格干预，改变 $X$ 会改变 $Y$，则 $X$ 是 $Y$ 的原因之一。合格干预 $I$ 应改变 $X$，并且除经由 $X$ 外不通过其他路径直接影响 $Y$；它还应切断 $X$ 与其通常原因的联系。

这种要求可以用图表示：

\begin{center}
\begin{tikzpicture}[node distance=14mm and 18mm]
\node[causalnode] (i) {$I$};
\node[causalnode,right=of i] (x) {$X$};
\node[causalnode,right=of x] (y) {$Y$};
\node[causalnode,above=of x] (u) {$U$};
\draw[causalarrow] (i)--(x);
\draw[causalarrow] (x)--(y);
\draw[causalarrow] (u)--(x);
\draw[causalarrow] (u)--(y);
\end{tikzpicture}
\end{center}

理想干预使 $I$ 只通过 $X$ 影响 $Y$，并替代 $U\rightarrow X$ 的通常生成方式。随机试验是实现这种结构的一种设计，但干预主义也评价不可实际操作的变量：只要能够清楚定义合适的反事实改变，例如地壳运动或恒星质量，就可以讨论其因果影响。
Woodward还以\term{不变性}说明解释深度。一条关系若在多种针对输入的干预下保持稳定，就比只在狭窄数据范围内拟合良好的相关式更具有因果解释力。温度计读数与温度稳定共变，但直接改变温度计指针不会改变环境温度，显示关系方向不是由回归对称性决定。
干预定义会否循环？要判断 $I$ 不通过其他路径影响 $Y$，似乎已经使用因果概念。Woodward的回应不是把所有因果词还原为非因果词，而是提供相互约束的网络分析：变量关系、干预条件与不变性共同澄清科学实践中的因果主张。这是一种非还原分析，而非传统意义上的概念消除。

\subsection{理论互补、解释层级与选择准则}

以药物安全为例，概率数据说明药物如何改变不良事件风险；潜在结果或图模型说明在何种设计下该变化可归因于药物；机制研究追踪代谢与受体通路；能力语言描述药物在不同背景中的倾向；实际因果分析判断某病例的不良事件是否由药物造成。这些答案不能相互代替，却可以形成证据链。
理论多元的纪律在于明确接口：机制知识如何转化为图结构假设？概率效应在哪些背景中保持？干预改变的是机制中的哪个环节？个例归因是否有群体数据与过程证据支持？只有回答这些接口问题，多种理论才是互补；否则“综合”只是把不同术语并列。

\paragraph{贯穿案例：医院内感染}

某医院发现中心静脉导管患者血流感染率升高。概率理论首先比较置管与未置管风险，但病情严重者更可能置管，简单概率提升包含适应证混杂。潜在结果与图模型要求定义“在符合置管条件者中采用某种导管策略”的效应，并调整病情等共同原因。
机制调查追踪皮肤菌群、置管过程、导管表面生物膜和免疫反应。它解释为何无菌操作、敷料和拔管时机可能改变风险。INUS分析可把“导管存在、无菌屏障破坏、足够留置时间、易感宿主”视为一条充分条件组合；另一条感染路径可能来自远处感染播散。没有哪个因素单独必要充分。
因果能力语言说导管表面具有支持生物膜的倾向，但这种能力在涂层、抗菌措施和宿主不同条件下表现不同。干预主义要求检查改变导管材料、操作规范或留置时间是否系统改变感染，并区分这些可控变量的路径。实际病例归因还要判断培养结果、时间与替代来源。

这个案例显示，证据不是简单投票。若观察关联强但机制检测显示主要病例菌株来自另一来源，因果解释需修订；若机制在实验室成立但医院随机改进措施无效果，可能是实施失败、机制贡献太小或替代路径补偿。不同证据之间的冲突用于诊断模型，而非选择自己偏好的理论。

\paragraph{解释层级与干预层级}

医院可以在分子、设备、操作人员、病房和制度层面干预。同一结果具有多层原因：材料属性影响附着，护士工作负荷影响无菌操作，采购制度影响材料选择。低层机制不一定使高层因果失真；若改变排班稳定改变感染，高层变量具有政策相关的干预意义。
层级之间可能出现实现关系。工作负荷通过多种具体动作影响感染，不存在一个固定微观中介；要求把政策效应完全还原为每次手部动作才算“真正因果”，会让有用的宏观解释失效。反过来，高层标签若无法对应任何稳定实施方式，效应也难迁移。
因果多元主义由此不是给每层一个互不相干真理，而是要求跨层约束：宏观干预改变哪些机制分布，微观机制为何汇总成观察效应，背景变化如何改变实现。机制、概率和干预关系在这些接口上互相检验。

\paragraph{理论选择的实用准则}

如果目标是政策净效果，优先明确处理策略与人群平均反事实；如果目标是故障诊断，优先追踪实际过程和部件；如果目标是风险预警，概率关系可以在不承担干预解释时发挥作用；如果目标是迁移，机制与不变性尤为关键；如果目标是责任，需在实际因果之外加入规范标准。
同一项目可随阶段改变工具：探索阶段以概率和发现算法缩小假设，验证阶段以试验识别差异，解释阶段分解机制，部署阶段检验环境不变性。理论的价值由它对具体查询的约束力衡量，而不是由它能否把所有因果语言压成一句定义衡量。

\section{理论选择、证据整合与模型责任}

\subsection{用反例检验理论}

因果理论的差异只有放入困难案例才会清楚。先占案例中，两条潜在路径都足以产生结果，但其中一条先完成过程；简单若无检验可能把被截断的备份也算原因。对称过度决定中，两项因素同时足以，删除任一项结果仍发生，若无条件又把两个直觉原因都排除。遗漏案例中，不行动可能是结果的原因，却难以用能量传递描述。共同原因案例则提醒概率提升可以由上游因素制造。

检验理论时应逐项记录它改变了什么。刘易斯式方案可能修改事件描述、引入因果链或采用影响关系；机制论可能要求追踪实际过程；结构模型可能加入应急情形和正常性排序；概率理论可能条件化背景集合。修正解决一个反例时，往往同时增加新的理论负担。例如加入正常性有助于区分园丁浇水与植物存活，却需要解释正常判断来自统计、功能还是规范。

反例不必迫使所有理论竞争一个定义。若研究任务是分配法律责任，实际因果和规范背景不可避免；若任务是估计药物平均效应，个例正常性可能不是主要对象；若任务是解释细胞过程，机制组织比总体风险差更重要。真正需要统一的是同一任务中的概念使用，而不是把所有因果问题压进一个必要充分条件。

\paragraph{多元证据如何形成共同约束}

理论多元主义容易滑向“每种说法都有道理”。受约束的多元主义则要求不同表示在共享部分相互校验。反事实模型给出若改变条件结果会怎样；机制模型说明该改变通过哪些活动传播；概率证据估计关系的总体强度与异质性；干预模型检查关系能否在外部改变下保持。若机制声称存在强而稳定的通路，干预数据却长期没有任何差异，就必须解释补偿、剂量不足或测量失败，不能各说各话。

整合的基本单位应是同一个查询。例如研究高温对老年死亡的影响，时间序列提供短期风险变化，自然实验利用异常天气，生理研究说明心血管与体温调节机制，建筑和社会资料刻画空调、居住隔离与照护网络。它们只有在暴露时间窗、目标人群和结局含义对齐后才能共同支持结论。若机制实验使用极端剂量，而人口研究关注日常波动，二者的表面一致不足以完成迁移。

证据之间也可分配否证任务。负对照针对某类未测混杂，过程测量针对路径时序，剂量反应针对机制强度，跨地区重复针对背景稳定性。没有单项检验能证明整个因果模型，但多项具有不同失败模式的证据可以缩小仍然可行的解释集合。多元主义的严格性体现在明确每项证据排除了什么，而非罗列文献数量。

\subsection{实在论、工具主义与模型责任}

因果实在论主张模型试图描述世界中真实的产生或依赖结构；工具主义更强调模型组织预测和行动的效用。实践中二者并非总是截然对立。一个用于政策的模型若在环境改变后持续失败，研究者会追问是否遗漏真实机制；一个声称刻画真实结构的模型，也必须展示它带来哪些可检验差异。争论焦点通常不是要不要接触现实，而是什么证据足以支持超出观测分布的结构承诺。

模型责任因此包含两面。表示责任要求说明节点、事件和机制为何以这种粒度划分，以及缺边和模块性假设凭什么成立。使用责任要求说明模型将服务预测、解释、干预还是归责，并防止输出跨越任务边界。一个为群体资源配置估计的平均效应，不应被直接解释为某个个体必然获益；一个用于机制探索的细胞模型，也不应未经迁移就支持全国政策。

采取这种立场不要求先解决全部形而上学争论。研究者可以承认因果结构有实在论意图，同时对当前模型保持可错性；也可以重视行动成功，却不把短期预测成功当作结构真理。关键是让结构承诺、经验风险和使用后果同时可审查。

\begin{chapterreview}
反事实理论把原因理解为在适当替代世界中制造结果差异的因素。Lewis用可能世界相似性解释反事实，以依赖链处理间接原因，并以“影响”分析捕捉结果发生方式的变化。先占、过度决定、遗漏和双重防止显示，简单若无检验不足以覆盖实际因果；事件描述、过程和正常性也可能进入判断。 概率理论适合处理不确定结果，但必须控制背景与混杂；INUS分析揭示复合充分路径；因果过程与机制论说明现象如何产生；力量论强调倾向在复杂背景中的表现；干预主义用合格变化和不变性连接因果与实验。它们聚焦不同查询，可在明确接口下共同构成因果证据。
\end{chapterreview}

\begin{selfcheck}
\begin{enumerate}
  \item 为什么材料条件句不能表达反事实？“最接近世界”需要保留哪些事实？
  \item 分别画出早期先占与对称过度决定的事件结构，说明若无检验为何失败。
  \item “作为影响的因果”如何利用结果发生方式，而不仅是结果是否发生？
  \item 选择一个遗漏因果案例，说明规范期待在判断中发挥什么作用。
  \item 构造一个 $P(E\mid C)>P(E\mid\neg C)$ 但 $C$ 不是 $E$ 原因的例子。
  \item 把一个现实故障写成至少两条充分路径，并识别其中的INUS条件。
  \item 为什么列出机制部件不等于给出机制解释？
  \item 对一个教育政策定义理想干预，并说明现实实施可能违反哪些条件。
\end{enumerate}
\end{selfcheck}

\chapterreadings{
反事实语义以\textcite{lewisCounterfactuals}为基础，早期因果分析见\textcite{lewisCausation1973}，后期修正见\textcite{lewisInfluence2000}。关于依赖与产生过程这两类因果概念的区分可读\textcite{hall2004}；实际因果的结构模型路线可在第11章后结合\textcite{halpern2016}阅读。

概率路线参阅\textcite{reichenbach1956}与\textcite{suppes1970}；复合条件分析见\textcite{mackie1974}；因果过程可对照\textcite{salmon1984}与\textcite{dowe2000}；新机制论可读\textcite{mdc2000}、\textcite{craver2007}与\textcite{glennan2017}；能力、行动者和干预路线分别以\textcite{cartwright1989}、\textcite{menziesprice1993}和\textcite{woodward2003}为核心。

当代争论的系统入口是\textcite{paulhall2013}，而\textcite{collinshallpaul2004}汇集了先占、过度决定和反事实语义的代表性分歧。对比式解释见\textcite{schaffer2005}，结构模型中的实际因果判准可从\textcite{halpernpearl2005}进入；解释为何有深浅之分则可读\textcite{strevens2008}。干预主义在生命科学层级问题中的适用条件见\textcite{woodward2010}。

机制路线不应停在定义：\textcite{craverdarden2013}展示机制如何在实际发现过程中被逐步约束，\textcite{glennanillari2018}提供跨科学领域的专题地图。力量论的系统方案见\textcite{mumfordanjum2011}；医学中的机制证据与统计证据怎样共同承担因果论证，可由\textcite{russowilliamson2007}作为批判性案例。
}
